\documentclass[11pt]{article}
\usepackage{amsmath,amssymb}
\usepackage[margin=1in]{geometry}
\usepackage{hyperref}
\emergencystretch=3em

\title{State density: integrability, Mellin transform and the leading coefficient (Stage 1c--d)}
\author{Claude, with Daniel Murfet}
\date{2026-09-07}

\begin{document}
\maketitle
\sloppy

\begin{abstract}
Units 225--226 complete Stage~1 of the Taylor-tree programme. Unit 225 exports the exact state
density $v=\mathrm{eval}(\mathrm{stateDensityRep}\,n\,w)$ of unit~224 in the forms later stages need:
integrability on $(0,1]$ (from the box identity with $g=1$, no analysis of individual power--log terms),
the real-valued identity $\int_{(0,1]^{n+1}}\prod a_i^{w_i}f(\prod a_i)\,da=\int_0^1v f$ for measurable
$f\ge0$ (both sides as \texttt{toReal} of the $\mathbb R_{\ge0}^\infty$ identity, no integrability side
conditions), the Mellin transform $\int_0^1z^sv(z)\,dz=\prod_i1/(w_i+s+1)$ (the paper's product zeta
function in the weight normalisation, by the same one-coordinate recursion), and the Mellin transform of
a basis term $\int_0^1\tau^{c-1}(-\log\tau)^j\,d\tau=j!/c^{j+1}$ (substitution $\tau=e^{-x}$ into Euler's
integral), hence termwise $\int_0^1z^s\,\mathrm{eval}\,c=\sum c_t\,j_t!/(s+\mu_t)^{j_t+1}$. Unit 226 proves
the \textbf{top coefficient}: with $l=\min(w_i+1)$ of multiplicity $m$,
\[
c_{l,m-1}=\frac1{(m-1)!}\prod_{w_i+1\ne l}\frac1{w_i+1-l},
\]
by an Abelian argument identical in spirit to Headline~XXI: $(s+l)^m\int_0^1z^sv$ tends, as $s\to-l^+$,
to that product (from the product transform) and to $(m-1)!\,c_{l,m-1}$ (termwise: every term but the
top one vanishes in the limit). In the monomial normalisation the coefficient is $a_{-m}/(m-1)!$,
$a_{-m}=\mathrm{mellinCoeff}\,h\,k\,l\,1$ (\texttt{monomial\_leadCoeff}): Astra~\#26's factorial
conversion between density and Laurent coefficients, as a theorem. Zero \texttt{sorry}/\texttt{axiom};
9103 jobs.
\end{abstract}

\section{The things formalised}
{\small
\begin{verbatim}
theorem weightedBoxIntegral_rpow : WBI d w (ofReal ∘ (·^s)) = ∏ ofReal (1/(wᵢ+s+1))
theorem stateDensity_integrableOn (hw : ∀ i, -1 < w i) : IntegrableOn (eval rep) (Ioc 0 1)
theorem integral_unitBox_eq_stateDensity (hf : Measurable f)
  (hf0 : ∀ z ∈ Ioc 0 1, 0 ≤ f z) :
  ∫ a in unitBox (n+1), (∏ aᵢ^{wᵢ}) * f (∏ aᵢ) = ∫ z in Ioc 0 1, eval rep z * f z
theorem mellin_stateDensity : ∫ z in Ioc 0 1, z^s * eval rep z = ∏ 1/(wᵢ+s+1)
theorem integral_Ioc_rpow_mul_neg_log_pow (hc : 0 < c) : ∫ τ^(c-1)(-log τ)^j = j!/c^(j+1)
theorem integrableOn_powLogBasis (hc : 0 < c) (j)
theorem mellin_eval (hc : ∀ t ∈ c, 0 < t.1 + s) : ∫ z^s eval c = ∑ cₜ jₜ!/(s+μₜ)^(jₜ+1)
def PowLogRep.coeffAt c μ j;  theorem term_tendsto; list_tendsto; product_tendsto
theorem stateDensityRep_leadCoeff (hl : ∀ i, l ≤ w i + 1) (hatt) :
  coeffAt (stateDensityRep n w) l (expMult w l - 1)
    = (1/(expMult w l - 1)!) * ∏ i, if w i + 1 = l then 1 else 1/(w i + 1 - l)
theorem monomial_leadCoeff : (∏ 1/(2kᵢ)) * coeffAt (rep (ratioExp − 1)) l (m−1)
    = mellinCoeff h k l 1 / (m−1)!
\end{verbatim}
}

\section{Mathematics}
The Abelian limit is the natural tool once the Mellin transform is a product: multiplying by
$(s+l)^m$ cancels the $m$ factors at the minimum, and the surviving product is continuous at $s=-l$
because the other exponents are strictly larger. On the density side the degree bound $j\le m-1$ at
$\mu=l$ (unit 224) is exactly what makes every lower-degree term vanish: $(s+l)^{m-1-j}\to0$ for
$j<m-1$. No inverse Mellin transform and no partial-fraction formula for the other coefficients are
needed for the leading term; the full residue formula (\texttt{eq:residuecoefficients}) remains a
corollary for later, checked numerically in the unit~224 replica.

\section{Lean notes}
\begin{itemize}
\item `\lambda` is the lambda token and cannot name a variable; use `l`.
\item \texttt{tendsto\_const\_nhds.div h hne} needs the constant fixed: \texttt{(tendsto\_const\_nhds (x := (1:ℝ))).div}.
\item \texttt{rw [show m = a + b by omega]} rewrites every \texttt{m}, including inside the exponent
  being manipulated; state the power identity as a separate \texttt{have} instead.
\item Real identity for nonnegative $f$ without integrability: \texttt{integral\_eq\_lintegral\_of\_nonneg\_ae}
  on both sides, then the $\mathbb R_{\ge0}^\infty$ theorem.
\item The substitution $\tau=e^{-x}$: \texttt{integral\_image\_eq\_integral\_abs\_deriv\_smul} on
  \texttt{Ioc 0 1} with \texttt{-log}, image \texttt{Ici 0}; its \texttt{iff} version transports integrability.
\end{itemize}

\section{Status}
Stage 1 complete (Headlines XXII, XXII$'$, XXII$''$). Next: review v18 of units 223--226; Stage 2 --- the
exact monomial moment identity with constant phase: $\int_{(0,1]^d}u^{h}(\sqrt n u^k)^p
e^{-\beta nu^{2k}+\beta\sqrt n u^k\xi_0}\,du=\sum_{\mu,j}\sum_{k'}n^{-\mu}(\log n)^{k'}c_{\mu,j}
\binom{j}{k'}\int_0^{n}t^{\mu+p/2-1}(-\log t)^{j-k'}e^{-\beta t+\beta\sqrt t\xi_0}\,dt$ exactly, then the
tail $\int_n^\infty$ exponentially small.
\end{document}
